\documentclass[10pt,a4paper,twoside]{ctexart}
\usepackage{xeCJK}
\usepackage{amsmath, amsthm}
\usepackage{listings,xcolor}
\usepackage{geometry}
\usepackage{fontspec}
\usepackage{graphicx}
\usepackage[colorlinks, linkcolor=black]{hyperref}
\usepackage{fancyhdr}
\usepackage{ctex}

\usepackage{relsize}
\renewcommand{\_}{\textscale{1}{\textunderscore}} %设置下划线长度

\setmainfont{Menlo}
\setmonofont{Menlo}
\setsansfont{Menlo}
\setCJKmainfont{Kai}

\ctexset{
    section={
        name={,、},
        number={\chinese{section}},
        format={\zihao{3}\centering},
        aftername={}
    }
}
\ctexset{
    subsection={
        format={\zihao{-4}},
        aftername={ }
    }
}
\ctexset{
    subsubsection={
        format={\zihao{-4}},
        aftername={ }
    }
}

\CTEXoptions[today=old]

\linespread{1.2}

\pagestyle{fancy}
% \lfoot{Xidian University}
% \rfoot{Leachim}
\fancyfoot[LE,RO]{\thepage}
\cfoot{}
\fancyfoot[RE]{Xidian University}
\fancyfoot[LO]{Leachim}

\renewcommand{\headrulewidth}{0.4pt} 
\renewcommand{\footrulewidth}{0.4pt}

\geometry{left=3.18cm,right=3.18cm,top=2.54cm,bottom=2.54cm} %A4打印页边距

\usepackage{textcomp}

\usepackage{minted}
\renewcommand{\theFancyVerbLine}{\sffamily \textcolor[RGB]{128,128,128}{\small\arabic{FancyVerbLine}}}

\setminted[c++]{
    breaklines=true,
    breakautoindent=false,
    breakanywhere=true,
    % escapeinside=\#\%,
    mathescape=true,
    linenos=true,
    tabsize=4,
    frame=single,
    style=arduino
}

\begin{document}
\title{\zihao{0}{ICPC Templates}}
\author {\zihao{2}{Leachim}\\\zihao{4}{Xidian University}}
\maketitle

\begin{figure}[htbp]
    \centering{\includegraphics[width=10cm]{../Figures/ICPC.png}}
\end{figure}

\newpage

\tableofcontents

\newpage
\section{图论}
\subsection{树的重心Get\_Centroid}
\inputminted{c++}{../图论/树的重心Get_Centroid.cpp}
\subsection{最小生成树Kruskal}
\inputminted{c++}{../图论/最小生成树Kruskal.cpp}
\subsection{最小生成树Prim}
\inputminted{c++}{../图论/最小生成树Prim.cpp}
\subsection{最短路Djikstra}
\inputminted{c++}{../图论/最短路Djikstra.cpp}
\subsection{最短路SPFA}
\inputminted{c++}{../图论/最短路SPFA.cpp}
\subsection{最近公共祖先LCA\_Doubling}
\inputminted{c++}{../图论/最近公共祖先LCA_Doubling.cpp}
\subsection{最近公共祖先LCA\_Tarjan}
\inputminted{c++}{../图论/最近公共祖先LCA_Tarjan.cpp}
\subsection{最近公共祖先LCA\_Heavy-Light\_Decomposition}
\inputminted{c++}{../图论/最近公共祖先LCA_Heavy-Light_Decomposition.cpp}
\subsection{判断负环SPFA\_Negtive\_Cycle}
\inputminted{c++}{../图论/判断负环SPFA_Negtive_Cycle.cpp}
\subsection{拓扑排序Topological\_Sort\_Khan}
\inputminted{c++}{../图论/拓扑排序Topological_Sort_Khan.cpp}
\subsection{树链剖分Heavy-Light\_Decomposition}
\inputminted{c++}{../图论/树链剖分Heavy-Light_Decomposition.cpp}

\newpage
\section{多项式}
\subsection{FFT字符串匹配String\_Match\_FFT}
\inputminted{c++}{../多项式/FFT字符串匹配String_Match_FFT.cpp}
\subsection{FFT递归Fast\_Fourier\_Transform\_Cooley-Tukey\_Recursion}
\inputminted{c++}{../多项式/FFT递归Fast_Fourier_Transform_Cooley-Tukey_Recursion.cpp}
\subsection{FFT递推Fast\_Fourier\_Transform\_Cooley-Tukey\_Iteration}
\inputminted{c++}{../多项式/FFT递推Fast_Fourier_Transform_Cooley-Tukey_Iteration.cpp}
\subsection{NTT递推Number\_Theoretic\_Transforms}
\inputminted{c++}{../多项式/NTT递推Number_Theoretic_Transforms.cpp}
\subsection{MTT任意模数Fast\_Fourier\_Transform\_MOD}
\inputminted{c++}{../多项式/MTT任意模数Fast_Fourier_Transform_MOD.cpp}
\subsection{多项式求逆Polynomial\_Inverse}
\inputminted{c++}{../多项式/多项式求逆Polynomial_Inverse.cpp}

\newpage
\section{字符串}
\subsection{字符串哈希String\_Hash}
\inputminted{c++}{../字符串/字符串哈希String_Hash.cpp}
\subsection{最小表示法Lexicographically\_Minimal\_String\_Rotation}
\inputminted{c++}{../字符串/最小表示法Lexicographically_Minimal_String_Rotation.cpp}
\subsection{马拉车Manacher}
\inputminted{c++}{../字符串/马拉车Manacher.cpp}
\subsection{字符串匹配KMP}
\inputminted{c++}{../字符串/字符串匹配KMP.cpp}
\subsection{AC自动机AC-Automaton}
\inputminted{c++}{../字符串/AC自动机AC-Automaton.cpp}
\subsection{后缀数组Suffix\_Array}
\inputminted{c++}{../字符串/后缀数组Suffix_Array.cpp}
\subsection{后缀自动机Suffix-Automaton}
\inputminted{c++}{../字符串/后缀自动机Suffix-Automaton.cpp}
\subsection{广义后缀自动机General\_Suffix-Automaton}
\inputminted{c++}{../字符串/广义后缀自动机General_Suffix-Automaton.cpp}

\newpage
\section{数据结构}
\subsection{并查集Disjoin\_Set\_Union}
\inputminted{c++}{../数据结构/并查集Disjoin_Set_Union.cpp}
\subsection{可撤销并查集Disjoin\_Set\_Union\_Withdrawable}
\inputminted{c++}{../数据结构/可撤销并查集Disjoin_Set_Union_Withdrawable.cpp}
\subsection{分块Block\_1}
\inputminted{c++}{../数据结构/分块Block_1.cpp}
\subsection{分块Block\_2}
\inputminted{c++}{../数据结构/分块Block_2.cpp}
\subsection{分块Block\_4}
\inputminted{c++}{../数据结构/分块Block_4.cpp}
\subsection{树状数组Binary\_Indexed\_Tree}
\inputminted{c++}{../数据结构/树状数组Binary_Indexed_Tree.cpp}
\subsection{二维树状数组2D\_Binary\_Indexed\_Tree}
\inputminted{c++}{../数据结构/二维树状数组2D_Binary_Indexed_Tree.cpp}
\subsection{线段树Segment\_Tree}
\inputminted{c++}{../数据结构/线段树Segment_Tree.cpp}
\subsection{线段树Segment\_Tree\_Multiply}
\inputminted{c++}{../数据结构/线段树Segment_Tree_Multiply.cpp}
\subsection{扫描线Scanline}
\inputminted{c++}{../数据结构/扫描线Scanline.cpp}
\subsection{zkw线段树ZKW\_Segment\_Tree}
\inputminted{c++}{../数据结构/zkw线段树ZKW_Segment_Tree.cpp}
\subsection{李超线段树Li-Chao\_Segment\_Tree}
\inputminted{c++}{../数据结构/李超线段树Li-Chao_Segment_Tree.cpp}
\subsection{可并堆左偏树Leftist\_Tree}
\inputminted{c++}{../数据结构/可并堆左偏树Leftist_Tree.cpp}
\subsection{Splay树Splay\_Tree}
\inputminted{c++}{../数据结构/Splay树Splay_Tree.cpp}
\subsection{Splay树Splay\_Tree\_Flip}
\inputminted{c++}{../数据结构/Splay树Splay_Tree_Flip.cpp}
\subsection{Splay树Splay\_Tree\_Dye\&Flip}
\inputminted{c++}{../数据结构/Splay树Splay_Tree_Dye\&Flip.cpp}
\subsection{动态树Link-Cut\_Tree}
\inputminted{c++}{../数据结构/动态树Link-Cut_Tree.cpp}
\subsection{红黑树Red\_Black\_Tree}
\inputminted{c++}{../数据结构/红黑树Red_Black_Tree.cpp}

\newpage
\section{数论}
\subsection{线性求逆元Modular\_Multiplicative\_Inverse}
\inputminted{c++}{../数论/线性求逆元Modular_Multiplicative_Inverse.cpp}
\subsection{数论分块Block\_Division}
\inputminted{c++}{../数论/数论分块Block_Division.cpp}
\subsection{贝祖引理Bezout\_Lemma}
\inputminted{c++}{../数论/贝祖引理Bezout_Lemma.cpp}
\subsection{卢卡斯Lucas}
\inputminted{c++}{../数论/卢卡斯Lucas.cpp}
\subsection{拓展欧几里得Exgcd}
\inputminted{c++}{../数论/拓展欧几里得Exgcd.cpp}
\subsection{拓展欧拉定理Ex\_Euler\_Theorem-Automaton}
\inputminted{c++}{../数论/拓展欧拉定理Ex_Euler_Theorem.cpp}
\subsection{中国剩余定理Chinese\_Remainder\_Theorem}
\inputminted{c++}{../数论/中国剩余定理Chinese_Remainder_Theorem.cpp}
\subsection{拓展中国剩余定理Ex\_Chinese\_Remainder\_Theorem}
\inputminted{c++}{../数论/拓展中国剩余定理Ex_Chinese_Remainder_Theorem.cpp}
\subsection{欧拉筛Eular\_Sieve}
\inputminted{c++}{../数论/欧拉筛Eular_Sieve.cpp}
\subsection{杜教筛Dujiao\_Sieve}
\inputminted{c++}{../数论/杜教筛Dujiao_Sieve.cpp}
\subsection{求原根Get\_Primitive\_Root}
\inputminted{c++}{../数论/求原根Get_Primitive_Root.cpp}
\subsection{素数测试Miller\_Rabin}
\inputminted{c++}{../数论/素数测试Miller_Rabin.cpp}
\subsection{大数分解Pollard\_Rho}
\inputminted{c++}{../数论/大数分解Pollard_Rho.cpp}


\newpage
\section{组合数学}
\subsection{康托展开Cantor}
\inputminted{c++}{../组合数学/康托展开Cantor.cpp}
\subsection{波利亚P\'{o}lya}
\inputminted{c++}{../组合数学/波利亚Polya.cpp}
\subsection{卡特兰数Catalan}
\begin{enumerate}
\item 有 $2n$ 个人排成一行进入剧场。入场费 5 元。其中只有 $n$ 个人有一张 5 元钞票，另外 $n$ 人只有 10 元钞票，剧院无其它钞票，问有多少中方法使得只要有 10 元的人买票，售票处就有 5 元的钞票找零？
\item 一位大城市的律师在她住所以北 $n$ 个街区和以东 $n$ 个街区处工作。每天她走 $2n$ 个街区去上班。如果他从不穿越（但可以碰到）从家到办公室的对角线，那么有多少条可能的道路？
\item 在圆上选择 $2n$ 个点，将这些点成对连接起来使得所得到的 $n$ 条线段不相交的方法数？
\item 对角线不相交的情况下，将一个凸多边形区域分成三角形区域的方法数？
\item 一个栈（无穷大）的进栈序列为 $1,2,3, \cdots ,n$ 有多少个不同的出栈序列？
\item $n$ 个结点可构造多少个不同的二叉树？
\item $n$ 个不同的数依次进栈，求不同的出栈结果的种数？
\item $n$ 个 $+1$ 和 $n$ 个 $-1$ 构成 $2n$ 项 $a_1,a_2, \cdots ,a_{2n}$ ，其部分和满足 $a_1+a_2+ \cdots +a_k \geq 0(k=1,2,3, \cdots ,2n)$ 对与 $n$ 该数列为？
\end{enumerate}

\begin{table}[ht]
\begin{tabular}{|l|l|l|l|l|l|l|}
\hline
H0 & H1 & H2 & H3 & H4 & H5 & H6  \\ \hline
1  & 1  & 2  & 5  & 14 & 42 & 132 \\ \hline
\end{tabular}
\end{table}


关于 Catalan 数的常见公式：

$$
H_n = \begin{cases}
    \sum_{i=1}^{n} H_{i-1} H_{n-i} & n \geq 2, n \in \mathbf{N_{+}}\\
    1 & n = 0, 1
\end{cases}
$$

$$
H_n = \frac{H_{n-1} (4n-2)}{n+1}
$$

$$
H_n = \binom{2n}{n} - \binom{2n}{n-1}
$$

$$
H_n = \frac{\binom{2n}{n}}{n+1}(n \geq 2, n \in \mathbf{N_{+}})
$$
\subsection{斯特林数Stirling}
\subsubsection{第一类斯特林数}

（斯特林轮换数）

$\begin{bmatrix}n\\ k\end{bmatrix}$ 表示将 $n$ 个两两不同的元素，划分为 $k$ 个非空圆排列的方案数。

递推式

$$
\begin{bmatrix}n\\ k\end{bmatrix}=\begin{bmatrix}n-1\\ k-1\end{bmatrix}+(n-1)\begin{bmatrix}n-1\\ k\end{bmatrix}
$$

边界是 $\begin{bmatrix}n\\ 0\end{bmatrix}=[n=0]$ 。

\subsubsection{第二类斯特林数}

（斯特林子集数）

$\begin{Bmatrix}n\\ k\end{Bmatrix}$ 表示将 $n$ 个两两不同的元素，划分为 $k$ 个非空子集的方案数。

递推式

$$
\begin{Bmatrix}n\\ k\end{Bmatrix}=\begin{Bmatrix}n-1\\ k-1\end{Bmatrix}+k\begin{Bmatrix}n-1\\ k\end{Bmatrix}
$$

边界是 $\begin{Bmatrix}n\\ 0\end{Bmatrix}=[n=0]$ 。

\subsubsection{上升幂与普通幂的相互转化}

我们记上升阶乘幂 $x^{\overline{n}}=\prod_{k=0}^{n-1} (x+k)$ 。

则可以利用下面的恒等式将上升幂转化为普通幂：

$$
x^{\overline{n}}=\sum_{k} \begin{bmatrix}n\\ k\end{bmatrix} x^k
$$

如果将普通幂转化为上升幂，则有下面的恒等式：

$$
x^n=\sum_{k} \begin{Bmatrix}n\\ k\end{Bmatrix} (-1)^{n-k} x^{\overline{k}}
$$

\subsubsection{下降幂与普通幂的相互转化}

我们记下降阶乘幂 $x^{\underline{n}}=\dfrac{x!}{(x-n)!}=\prod_{k=0}^{n-1} (x-k)$ 。

则可以利用下面的恒等式将普通幂转化为下降幂：

$$
x^n=\sum_{k} \begin{Bmatrix}n\\ k\end{Bmatrix} x^{\underline{k}}
$$

如果将下降幂转化为普通幂，则有下面的恒等式：

$$
x^{\underline{n}}=\sum_{k} \begin{bmatrix}n\\ k\end{bmatrix} (-1)^{n-k} x^k
$$
\subsection{范德蒙德卷积Vandermonde\_Convolution}
$$
\sum\limits_{i=0}^k\binom{n}{i}\binom{m}{k-i}=\binom{n+m}{k}
$$

从数量分别为$n$和$m$的石堆中总共选取$k$个石子。


\newpage
\section{网络流}
\subsection{最大费用流Minimum-Cost\_Flow\_Edmonds-Karp}
\inputminted{c++}{../网络流/最大费用流Minimum-Cost_Flow_Edmonds-Karp.cpp}
\subsection{最大流Maximum\_Flow\_Edmonds-Karp}
\inputminted{c++}{../网络流/最大流Maximum_Flow_Edmonds-Karp.cpp}
\subsection{最大流Maximum\_Flow\_Dinic}
\inputminted{c++}{../网络流/最大流Maximum_Flow_Dinic.cpp}
\subsection{二分图最大匹配Bipartite\_Graph\_Maximum\_Matching\_Dinic}
\inputminted{c++}{../网络流/二分图最大匹配Bipartite_Graph_Maximum_Matching_Dinic.cpp}
\subsection{二分图最大匹配Bipartite\_Graph\_Maximum\_Matching\_Hungarian}
\inputminted{c++}{../网络流/二分图最大匹配Bipartite_Graph_Maximum_Matching_Hungarian.cpp}
\subsection{二分图最大匹配Bipartite\_Graph\_Maximum\_Matching\_Hopcroft-Karp}
\inputminted{c++}{../网络流/二分图最大匹配Bipartite_Graph_Maximum_Matching_Hopcroft-Karp.cpp}

\newpage
\section{计算几何}
\subsection{计算几何Compuational\_Geometry}
\inputminted{c++}{../计算几何/计算几何Compuational_Geometry.cpp}
\subsection{圆的K次面积交K-Intersection\_Area\_of\_Circles}
\inputminted{c++}{../计算几何/圆的K次面积交K-Intersection_Area_of_Circles.cpp}

\newpage
\section{优化算法}
\subsection{wqs二分+决策单调性dp}
\inputminted{c++}{../优化算法/wqs二分+决策单调性dp.cpp}
\subsection{斜率优化Slope\_Optimization}
\inputminted{c++}{../优化算法/斜率优化Slope_Optimization.cpp}

\newpage
\section{离线算法}
\subsection{莫队算法Mo}
\inputminted{c++}{../离线算法/莫队算法Mo.cpp}
\subsection{线段树分治Segment\_Tree\_Partition}
\inputminted{c++}{../离线算法/线段树分治Segment_Tree_Partition.cpp}

\newpage
\section{杂项}
\subsection{快读快写Fast\_Read\&Write}
\inputminted{c++}{../杂项/快读快写Fast_Read\&Write.cpp}
\subsection{快速幂Fast\_Power}
\inputminted{c++}{../杂项/快速幂Fast_Power.cpp}
\subsection{矩阵快速幂Matrix\_Fast\_Power}
\inputminted{c++}{../杂项/矩阵快速幂Matrix_Fast_Power.cpp}
\subsection{矩阵加速Matrix\_Acceleration}
\inputminted{c++}{../杂项/矩阵加速Matrix_Acceleration.cpp}
\subsection{最长公共子序列Longest\_Increasing\_Subsequence}
\inputminted{c++}{../杂项/最长公共子序列Longest_Increasing_Subsequence.cpp}
\subsection{模拟退火Simulated\_Annealing}
\inputminted{c++}{../杂项/模拟退火Simulated_Annealing.cpp}
\subsection{快速沃尔什变换Fast\_Walsh\_Transform}
\inputminted{c++}{../杂项/快速沃尔什变换Fast_Walsh_Transform.cpp}
\subsection{快速莫比乌斯变换Fast\_Mobius\_Transform}
\inputminted{c++}{../杂项/快速莫比乌斯变换Fast_Mobius_Transform.cpp}
\subsection{快速子集变换Fast\_Subset\_Transform}
\inputminted{c++}{../杂项/快速子集变换Fast_Subset_Transform.cpp}
\subsection{24点24\_Point}
\inputminted{c++}{../杂项/24点24_Point.cpp}




\end{document}